We now sharpen the count in \textbf{Lemma~\ref{lem:multigraph-cycle-counts}}
for odd $D$. Fix an induced red five-cycle and use its sets $U_i$,
their union $B$, and the number $L$ of blue pairs between neighboring
parts as in that lemma. Write $\mu=\frac{5}{4}D^2-m$. Since $m=f(D)+1$ and $D$ is odd,
$2\mu=D-\frac{5}{2}$. We first bound the contribution to
\eqref{eq:multi-core-cut} from outside vertices. We then show that a cycle
with small $R$ would force $m\le f(D)$.

In the graph on $B$ consisting of the marked pairs, the minimum size
of a vertex cover is at least $D+q$, where $q=a_*-\frac{L}{a_*}$
when $a_*>0$. In fact, any independent set in this graph with $x_i$
vertices in $U_i$ satisfies $\sum_i x_i x_{i+1}\le L$, and
\[
 \sum_i x_i x_{i+1}\ge a_*\left(\sum_i x_i-D\right).
\]
The difference between the two sides is affine in each $x_i$ on
$[0,a_i]$, so it suffices to check the inequality for each $x_i$ in
$\{0,a_i\}$. At such a point, delete one occupied class from each
consecutive occupied pair until the occupied classes are pairwise
nonconsecutive. An edge with both classes occupied accounts for a
product $a_i a_{i+1}$, and either deletion removes at most
$\frac{a_i a_{i+1}}{a_*}$ vertices. The remaining classes contain at
most $D$ vertices because any two nonconsecutive classes have a
common neighboring class. Hence any independent set has size at most
$D+\frac{L}{a_*}$. Subtracting from $|B|=2D+a_*$ gives the bound on
the vertex cover number. The argument applies also with enlarged
classes below. The core neighborhoods of the two ends of an outside
marked edge cover all marked pairs in $B$, so $Q_x+Q_y\ge D+q$ for
every such edge $xy$.

For a sharper estimate, let $q\ge0$ be any lower bound on the vertex
cover number minus $D$. Vertices of signature size zero, and vertices
of signature size one with no marked neighbor in $B$ or of signature
size two, now send only their positive contributions. They distribute
these contributions uniformly over their marked copies, all of which
meet vertices of signature size at least three. At a sender $x$ with
positive contribution, $b_x=Q_x$, so $Q_x<h_x\le D-Q_x$.
Each of its marked neighbors has at least $D+q-Q_x$ core neighbors
and hence at most $Q_x-q$ outside marked copies. In particular,
$q\le Q_x<\frac{D}{2}$. If $q\ge\frac{D}{2}$, there is no positive
transfer. Set $u=(\frac{D}{2}-q)_+$. A recipient receives at most
\[
 C(u)=\frac{2u^2}{\frac{D}{2}+u}.
\]
Indeed, the amount sent per copy is
$1-\frac{Q_x}{h_x}\le\frac{D-2Q_x}{D-Q_x}\le\frac{D-2q}{D-q}$,
and the recipient has at most $u$ outside marked copies. If $a_*>0$
and $q\ge a_*-\frac{L}{a_*}$, then $u\le\frac{R}{2}+\frac{L}{a_*}$.
Since $0\le C'(u)\le2$ for $u\ge0$, it follows that
$C(u)\le\frac{R^2}{D+R}+\frac{2L}{a_*}$.
More generally, if there are $c$ possible recipients, positive
transfer requires $Q_x<h_x\le c(Q_x-q)$. This is impossible for $c=1$,
and for $c=2$ it implies $Q_x>2q$. For two recipients the amount
sent per copy is at most $\frac{D-4q}{D-2q}$ whenever transfer
is possible. Multiplying by $u$ shows that each recipient receives
at most $\max\{0,\frac{D-4q}{2}\}$.

For each $x$ outside $B$, discard any negative contribution left
at a sender and denote by $\ell_x$ the difference between
$\frac{k_xD}{2}$ and the resulting contribution of $x$.
A vertex of signature size one has $\ell_x\ge t+2$ for
$t=\frac{D}{2}-R$ and $R\ge3$. Indeed, it either retains at most
$R-2$ or sends its whole positive contribution and retains zero.
For signature size two, $\ell_x=D-h_x+b_x\ge2b_x\ge2$.
A recipient of signature size $k$ has
$\ell_x\ge(\frac{k}{2}-1)D-C(u)$, or the bound obtained using the
smaller transfer estimate when $c\le2$. Substituting
$\sum_x(h_x-b_x)\le\frac{DK}{2}-\sum_x\ell_x$ into
\eqref{eq:multi-core-cut} and using $|B|=\frac{5D-R}{2}$ gives
\begin{equation}\label{eq:multi-loss-budget}
 2L+\frac{D}{2}(R-K)+\sum_x\ell_x\le2\mu.
\end{equation}

We enlarge the five sets $U_i$ while retaining all adjacencies between neighboring sets.

Every vertex $x$ whose signature consists of two nonconsecutive indices $i,j$,
with $b_x<a_*$, and with complete host adjacency to both $U_i$ and
$U_j$, is put into their common neighboring class. We show that
every pair joining consecutive enlarged classes is still an edge
of the simple support. For two newly
added vertices $x,y$ in consecutive classes $p,p+1$, choose a marked
neighbor of $x$ in $U_{p-1}$ avoiding the at most $s_{p+1}$ neighbors
of $y$ there. Then choose a marked neighbor of $y$ in $U_{p+2}$ avoiding
the neighbors of the first choice and of $x$. The latter two vertices
have at most $s_{p-1}+s_p$ neighbors in that class. Both choices
are possible by \eqref{eq:multi-core-identities} and $b_x,b_y<a_*$.
The two chosen marked edges force $xy$. The old--new pairs are
complete by the inclusion rule. A prime denotes enlarged quantities.
In particular each inclusion reduces $R$ and $K$ by two. The old
marked anchor stars can still be used.

Here are the comparisons for a remaining vertex of signature size two. If its
positions $\{i,j\}$ are nonconsecutive, $b_x<a_*$, and
it remains outside, it omits at least one point of both $U_i$ and $U_j$.
Indeed an omitted point in $U_i$, through its marked spoke, bounds the
number of distinct marked neighbors of $x$ in $U_j$ by $s_i$. This
contradicts complete adjacency to $U_j$ since $a_j-b_x>s_i$.
The same comparison with an omitted point works for neighbors in
$U'_j$, using the new bound $s_i'$. Its marked core neighbors can only
lie in $U_i'\cup U_j'$. To see this, take positions $\{4,1\}$, so that
comparison with the cycle excludes $U_0'$. A mark to $U_2'$ would force
at least $a_4-s_2'$ host neighbors of $x$ in $U_4$ by the star from
$c_0$. At most $s_1'$ of these are marked, contrary to $b_x<a_*$
by \eqref{eq:multi-core-identities}. A mark to $U_3'$ is excluded by
reflection. Both anchor pairs of $x$ are blue. If one were marked,
the marked star into the other designated class, based at an anchor
missing the first, would force complete adjacency to that other
class. Each of the two marked-neighbor counts into the designated
classes is at most the slack of the opposite class. Thus
\begin{equation}\label{eq:multi-low-blue-outside}
 Q_x\le s_i'+s_j'+b_x.
\end{equation}
If $xy$ is marked with $y\in U_i'$, every omitted old $U_j$ vertex must
see $y$, again by its marked spoke. At most $s_i'-1$ such vertices
can be omitted since $xy$ already consumes one unit at $y$. It follows
that $b_x\ge a_j-2s_i'+1\ge a_*-s_i'+1$. The same omitted vertices
also give
\begin{equation}\label{eq:multi-rectangle}
 I+b_{x,j}\ge a_j-s_i'\ge a_*,
\end{equation}
where $I$ counts host edges inside the enlarged core other than its
reserved base edges and
$b_{x,j}$ counts the blue neighbors of $x$ in old $U_j$.

A vertex with $b_x<a_*$ and signature $\{0,1\}$ can have marked edges
into the enlarged core only to $U_3'$. The cycle excludes marked edges to $U_0'$ and
$U_1'$. A mark to $U_2'$ forces at least $a_0-s_2'$ blue neighbors
in $U_0$ by the star from $c_4$, and a mark to $U_4'$ is excluded
by the reflected comparison. If there is a marked neighbor in $U_3'$,
the two ends of that mark must cover the marked $U_0$--$U_1$ pairs.
If the union of their neighborhoods omits $t_0$ vertices of $U_0$
and $t_1$ vertices of $U_1$, then all $t_0t_1$ pairs between these
omitted sets are blue. Thus $t_0t_1\le L$ and one of $t_0,t_1$ is at
most $\sqrt L$.
The neighbor in $U_3'$ can pay for at most $s_3'-1$ wrong-class
adjacencies in either part after accounting for the selected mark.
Thus $b_x\ge a_*+1-\sqrt L$. Finally, if a remaining outside vertex
has no enlarged core mark, a marked outside neighbor has $Q\le D-1$.
The old vertex-cover bound therefore gives $b_x=Q_x\ge q+1$.

Reserve one base copy for each pair in consecutive enlarged
classes, choosing a marked copy whenever possible. Let $L'$
count the unmarked base copies. Let $I$ count the other host
edges inside the enlarged core, $P$ count the marked ones among
them, $L_{\rm cut}$ count the unmarked copies across the cut of
the enlarged core, and $F_{\rm out}$ count the marked copies with both ends
outside. There are $\frac{D|B'|-A'}{2}-L'$ marked base copies.
The number of marked copies across the cut is at most
$A'-2I-L_{\rm cut}$, since the edges counted by $I$ each use two
units of the remaining core capacity $A'$. Adding $P$ and $F_{\rm out}$ and
using \eqref{eq:multi-core-identities} gives
\begin{equation}\label{eq:multi-extra-capacity}
 \mu\ge \frac{R'D-2A'}{4}+L'+(2I-P)+L_{\rm cut}-F_{\rm out}
 \ge L'+I+L_{\rm cut}-F_{\rm out} .
\end{equation}
When $R'=1$, all but one class have zero extra capacity. If all
outside signatures are singletons, outside marks are impossible.
Any non-base mark in the only positive-capacity class or from
outside to that class misses the opposite cycle edge, so only marked
base copies remain. Their total is at most $f(D)$.

\begin{lemma}\label{lem:half-degree-cycle}
If $D\ge9$ is odd and $m=f(D)+1$, then $R>\frac{D+1}{2}$ on
every induced five-cycle of red edges in the simple support.
\end{lemma}

\begin{proof}
Suppose that $R\le\frac{D+1}{2}$. Then $t=\frac{D}{2}-R\ge-\frac{1}{2}$.
If $R=1$, outside signatures have size at most
one and no outside vertex can have signature zero in the marked
endpoint host. The $R'=1$ base count above applies without enlargement.
Assume $R\ge3$ and $t\ge-\frac{1}{2}$. If $t\ge\frac{1}{2}$,
\eqref{eq:multi-rough-discharge} gives $L\le \frac{D}{2}-\frac{5}{4}$ and $R-K\le1$.
Also $a_*\ge3$ and $\frac{L}{a_*}<2$. For $t=-\frac{1}{2}$ and $D\ge13$, it gives
$L\le\frac{5D-9}{8}$ and $\frac{L}{a_*}<\frac{5}{2}$, with $a_*\ge3$. In both cases
$q>0$. The remaining case $(D,R)=(9,5)$ has $a_*=2$, $L\le4$ and
$q\ge0$. Also $K\le5$, so at most one recipient is possible there.
In all cases $C(u)\le \frac{D}{2}$ and the losses in
\eqref{eq:multi-loss-budget} are nonnegative. That inequality gives
$R-K\in\{0,1\}$ and $L\le \frac{D}{2}-\frac{5}{4}$.

There is at most one vertex of signature size at least three.
Indeed, $\frac{R^2}{D+R}\le \frac{D}{6}$ for $t\ge\frac{1}{2}$ and is at most $\frac{D}{6}+\frac{1}{2}$
for $t=-\frac{1}{2}$. The loss of three recipients, with $2L$ included,
would be at least $D-\frac{3}{2}$, since $a_*\ge3$. With two recipients,
the improved bound gives total loss with $2L$ at least $D-1$.
Both bounds exceed the budget. With one recipient no positive
transfer occurs, and its size is exactly three.

Suppose first there is one vertex of signature size three, say $z$. Let $a$
count vertices of signature size one. Parity and the budget force $a$ even
and $R-K=0$. The lower bound on the total loss is
$\frac{D}{2}+R-3+a(t+1)$. For $a\ge2$ this exceeds $D-\frac{5}{2}$, so $a=0$.
We show that every vertex with $k_x=2$ is included in the enlarged core.
For $t=-\frac{1}{2}$ the budget forces $L=0$ and each vertex of signature size two
to have one blue neighbor. A remaining vertex with nonconsecutive
signature positions and $b_x<a_*$ would have two blue anchor pairs.
The other remaining types
violate $b_x\ge a_*+1-\sqrt L$ or $b_x\ge q+1$. For $t\ge\frac{1}{2}$ a
remaining vertex with nonconsecutive signature positions, $b_x<a_*$
and a marked edge to the core satisfies
$b_x\ge a_*+1-s_i'\ge t+\frac{3}{2}$. Indeed, $R'=K'$ and the slack
at each position counts the remaining vertices that use it.
There are $\frac{R'-3}{2}$ such vertices of signature size two,
and at most one more can use that position, namely $z$.
Together with the loss of $z$ and the other vertices of signature size two
this gives a loss of at least $D+t-2>D-\frac{5}{2}$. A remaining vertex with $b_x\ge a_*$
costs at least $2a_*$ and a remaining vertex of either other type
costs at least $2a_*+2$ with $2L$ included. These also exceed the
budget after adding the other vertex losses.

The resulting enlarged core has $R'=3$, with one unit at each of the
three positions of $z$. Up to rotation and reflection, the two omitted
anchors are adjacent or nonadjacent, as shown in
\textbf{Figure~\ref{fig:cycle-three-neighbors}} of the \emph{Supplementary Proofs}.
Write $D=2d+1$. If the omitted positions are
nonconsecutive, the three classes of positive capacity have sizes
$d+1,d,d$, base count $f(D)-D$ and extra capacity $3d+1$.
If $zy$ is marked outside, then the vertex of signature size zero $y$ has
only blue core pairs. The two omitted anchor stars force
$Q_z'+Q_y'\ge3d+1-L'$, since each omitted point of those three
classes consumes a distinct blue base pair. Also $F_{\rm out}\le D-Q_z'$.
Thus $L'+L_{\rm cut}-F_{\rm out}\ge d$. If there is no outside mark, only
$z$ remains outside. The extra marks either all touch $z$ and
number at most $D$, or there is a marked non-base copy inside
the core. No same-class mark is possible by the opposite cycle
edge. A mark from the class of size $d+1$ to one of size $d$
blocks any mark from $z$ to those two classes. Its endpoints use
their only extra units and every possible edge joining their marked edge to
such a mark from $z$ would use one of those units. The total number
of extra marks is then at most $\frac{3d+1+d}{2}\le D$. If there are
only repeated internal marks between the two consecutive size-$d$
classes, one of those repetitions blocks marks from $z$ to the
size-$(d+1)$ class by the same comparison. The extra marks then
use at most $2d$ units and still number at most $D$.

If the two omitted positions are consecutive, put the three
positive positions at $0,1,2$. The base count is $f(D)-(D-1)$.
No outside mark sees the omitted cycle edge. A mark from $z$ to
$U_1'$ would also miss that edge, so $h_z\le D-1$ and $z$ marks
to $U_0'$ or $U_2'$. A mark to $U_0'$ forces $z$ to see every
vertex of $U_2'$ having a marked neighbor in $U_3'$ by comparison
with the marked edges to $U_3'$, where $s'_3=0$. If such a
neighbor is absent at some vertex of $U_2'$ then $L'\ge d+1$
and \eqref{eq:multi-extra-capacity} suffices. Otherwise the
adjacency to $U_2'$ is complete. If these pairs are all blue the
cut costs at least $d$. If one is marked, reflection forces
complete adjacency to $U_0'$ too. The two classes now have no
remaining extra capacity and the middle class cannot have a
same-class mark. The only extra marks meet $z$ and number at most
$D-1$. In every subcase $\mu\ge d$ suffices, since its fractional
part is $\frac{1}{4}$.

Suppose now that there is no vertex with three or more positions.
Then there are no vertices of signature size zero. If all outside vertices have
two positions, the budget and parity give $R-K=1$ and there are
$n=\frac{R-1}{2}$ outside vertices. The total number of unmarked host
edges is at most $\frac{D-5}{4}$. Indeed, summing the differences
between $D$ and the marked degree over the $\frac{5D-1}{2}$ vertices
counts at least twice the number of unmarked host edges.
Hence $\sum b_x\le \frac{D-5}{4}-L<a_*-\sqrt L$,
using $a_*\ge\frac{D-1}{4}$. Every outside vertex has a marked core
neighbor by the vertex-cover bound, while the comparison for
consecutive positions excludes those positions. Apply the enlargement
rule to every eligible vertex. Two remaining vertices cannot be joined
by a marked edge.
Their combined old-core neighborhood would have to omit at most
$\sqrt L$ points in at least three classes, since two consecutive
omitted sets have product at most $L$. In any class outside the
intersection of their two designated class sets, their combined
neighborhood omits at least $a_*-b_x-b_y>\sqrt L$ points. The
intersection has at most two classes. If no vertex remains, $R'=1$.
Otherwise take a marked edge to $U_i'$ from a remaining vertex $x$
of positions $\{i,j\}$. Each original vertex of signature size two pays
one blue anchor pair to $L'+L_{\rm cut}$. For $x$ choose its pair
to $c_i$, which is disjoint from its $b_{x,j}$ pairs. By
\eqref{eq:multi-rectangle} and \eqref{eq:multi-extra-capacity},
$\mu\ge n+a_*=\frac{D-1}{2}$, as required.

It remains to consider $a>0$ vertices of signature size one and no
vertex of signature size three. If $a$ is even then $R-K=1$ and the loss
is at least $\frac{D}{2}+R-1+a(t+1)>D-\frac{5}{2}$. Hence $a$ is odd and $R=K$.
Apply the enlargement rule and let $n$ be the number of remaining
vertices of signature size two, so $R'=a+2n$. After subtracting the base
losses, \eqref{eq:multi-loss-budget} implies
\[
 2L+\sum_{\rm remaining\ two}(\ell_x-2)
 \le T=2a_*-\frac{5}{2}-a(t+1).
\]
A vertex with $b_x\ge a_*$ contributes at least $2a_*-2>T$ to the
left side. For a vertex with $b_x<a_*$ and consecutive signature
positions, or a vertex without a marked edge to the core, its
contribution together with $2L$ is at least $2a_*$ by the bounds
$a_*+1-\sqrt L$ and $q+1$, respectively. Therefore each remaining
vertex with $k_x=2$ has $b_x<a_*$, nonconsecutive signature positions,
and a marked edge to the core. Its two anchor pairs are blue.

If two such vertices $x,y$ mark to one another, their signatures
are distinct. Equation \eqref{eq:multi-low-blue-outside} shows that
their marked-neighbor counts in the old core sum to at most
$2R'-n$, since any shared position has $s_i'\le a+n$.
Consequently
\[
 b_x+b_y\ge D+q-2R'+n\ge a_*+2t+n-\frac{L}{a_*}.
\]
The budget gives $b_x+b_y\le \frac{T}{2}-n+4-L$. These contradict for
$n\ge2$, $a\ge1$ and $t\ge-\frac{1}{2}$.

Suppose $xy$ is marked, where $k_x=1$ and $y$ is a remaining
vertex with $k_y=2$. After a rotation of the cycle, the signature
of $x$ is $\{0\}$ and that of $y$ is $\{1,3\}$ or $\{2,4\}$. The omitted anchor stars require
complete combined adjacency to the three indicated old classes.
If $x$ has a marked edge to $B$, it has at least $a_*+1$ blue pairs
in $U_0$. The designated class $U_1$ or $U_4$ for $y$ is short of
at least $a_*-b_y$ of its vertices. These must be covered by
additional blue pairs of $x$, since $x$ only marks to $U_2$ or
$U_3$ in the old core. If $x$ has no marked edge to $B$, the
three required classes have at least $D+a_*$ vertices and $y$
has at most $R'$ marked neighbors there. Both cases imply
$b_x+b_y\ge2a_*+1$. The budget then requires at least
$D+a(t+1)-2>D-\frac{5}{2}$ in losses after including the other vertices.
There are also no marks between two vertices of signature size one.
Thus $F_{\rm out}=0$. If $n>0$, then \eqref{eq:multi-rectangle} and the blue edges
from one vertex of signature size one to $B$ give $I+L_{\rm cut}\ge2a_*+1>\frac{D}{2}$.
If $n=0$ and $a\ge3$ then $L_{\rm cut}\ge3(a_*+1)>\frac{D}{2}$.
The last possibility is $R'=a=1$, already settled.


\end{proof}

\begin{lemma}\label{lem:odd-slack-refinements}
If $m=f(D)+1$, then each induced five-cycle of red edges satisfies
$R\ge9$ for $D=11$ and $R\ge11$ for $D\in\{13,15\}$.
\end{lemma}

\begin{proof}
By \textbf{Lemma~\ref{lem:half-degree-cycle}} and the odd parity of $R$,
it suffices to exclude the pairs $(11,7)$, $(13,9)$, and $(15,9)$ for $(D,R)$.
For $(11,7)$ and $(15,9)$, write $D=4k+3$ and $R=2k+3$.
Here $k=2,3$ and $a_*=k$. Equation \eqref{eq:multi-rough-discharge}
gives $L\le9,13$, respectively. Let $\tau$ be the minimum size
of a vertex cover of the marked pairs in $B$. If $L<D$, then
\begin{equation}\label{eq:multi-small-cover}
 \tau\ge D+a_*-\max\{0,\lfloor\sqrt L\rfloor-1\}.
\end{equation}
To prove it, consider an independent set of size $\alpha>D$ in
the marked old core, with $x_i$ points in $U_i$. It occupies a
consecutive pair of positions, since any two nonconsecutive classes
have total size at most $D$. If the occupied positions induce a
connected graph, summing $x_i x_j\ge x_i+x_j-1$ on a spanning tree
gives $L\ge\alpha-1\ge D$. If they induce a disconnected graph,
the occupied positions consist of one edge and one isolated point.
The two classes of the edge cannot contribute their own anchors,
because the anchor spokes to the adjacent class are marked. Either
of these classes and the isolated class now contribute at most
$D-1$ points. Hence the adjacent classes
each contribute at least $\alpha-D+1$ and $L\ge(\alpha-D+1)^2$.
Subtracting the upper bound for $\alpha$ from $|B|$ proves
\eqref{eq:multi-small-cover}.

Two further estimates apply to each of the three degree pairs. First, each absorbed vertex
$x$ contributes at most $D-h_x$ new blue consecutive pairs and at
least one such pair, since one of its two anchor pairs is blue.
Thus, if $p$ vertices are absorbed,
\begin{equation}\label{eq:multi-new-blue}
 L+p\le L'\le L+\sum_{x\ {\rm absorbed}}(D-h_x).
\end{equation}
The vertex cover number of the enlarged marked core is at least
$\tau'=\lceil D+a_*'-\frac{L'}{a_*'}\rceil$,
where $a_*'=\frac{D-R'}{2}$. Second, if all $F_{\rm out}$ outside marked
edges meet one vertex and that vertex has at most $p$
distinct outside marked neighbors, then
\begin{equation}\label{eq:multi-center-cut}
 F_{\rm out}\le\frac{p}{p+1}(2D-\tau').
\end{equation}
Indeed, one of its neighbors has outside marked degree at least
$\frac{F_{\rm out}}{p}$, and the outside marked degrees of the two ends sum to
at most $2D-\tau'$. The bound is immediate when $p=0$.

For $D=11,15$, the cover bound \eqref{eq:multi-small-cover} and the loss
budget \eqref{eq:multi-loss-budget} give $L\le2k$ and
$R-K\in\{0,1\}$. For $(D,R)=(13,9)$ the corresponding bounds are
$L\le7$ and $R-K\le2$. In each case $\sum_{x\notin B}k_x=K\le R$.
Thus there are only finitely many possible lists of positive signature
sizes. An outside vertex of signature size zero can have a marked
neighbor only among the vertices of signature size at least three.
The cover bound limits how many marked copies these neighbors can
receive. For the remaining vertices, \eqref{eq:multi-new-blue} bounds
the additional blue pairs after enlargement and \eqref{eq:multi-center-cut}
bounds the outside marks when they all meet one vertex.
\textbf{Lemma~\ref{lem:supp-slack-cases}} of the \emph{Supplementary Proofs}
carries out the resulting counts for the three pairs $(D,R)$ and
excludes each pair. This proves the claimed lower bounds on $R$.
\end{proof}
